\chapter{Modélisation mathématique}
\label{chap:modelisation}

Ce chapitre construit le modèle d'optimisation. La rédaction suit une discipline stricte :
\emph{rien n'est utilisé avant d'être défini}. On procède dans l'ordre ensembles,
paramètres, variables, contraintes, objectif. Chaque contrainte est accompagnée de sa
traduction en français, de l'organe industriel qu'elle représente, et de la justification de
sa forme.

\section{L'unité de compte du modèle}
\label{sec:unite}

Avant toute formalisation, il faut trancher une question que la documentation fournie
n'aborde nulle part : \emph{en quelle unité comptons-nous ?}

\subsection{L'énigme}

La documentation affirme que la concentration de \SI{29}{\percent} à \SI{54}{\percent}
s'effectue avec un rendement massique de \SI{100}{\percent}. C'est physiquement impossible
en masse d'acide : concentrer, c'est retirer de l'eau, donc perdre de la masse.

\subsection{La résolution : trois preuves convergentes}

\begin{proposition}
Toutes les grandeurs massiques du système sont exprimées en \textbf{tonnes de \PdeuxOcinq{}}
et non en tonnes d'acide marchand.
\end{proposition}

\paragraph{Preuve 1 --- par la structure des formules de conversion.}
La documentation fournit les conversions hauteur de cuve $\to$ tonnage :
\[
\mathrm{vol}_{29}(h) = 174{,}11 \cdot h \cdot 0{,}33
\qquad
\mathrm{vol}_{54}(h) = 95 \cdot h \cdot 0{,}83
\]
et, séparément, les paramètres physiques des cuves. Or
\[
\underbrace{0{,}26}_{\text{titre \PdeuxOcinq{}}} \times \underbrace{1{,}270}_{\text{densité}} = 0{,}3302 \approx 0{,}33
\qquad\text{et}\qquad
0{,}50 \times 1{,}660 = 0{,}8300 \;\text{(exact)}
\]
La structure de la conversion est donc
\[
\mathrm{vol}(h) = \underbrace{S \cdot h}_{\text{volume}} \times \underbrace{\rho}_{\text{densité}} \times \underbrace{\theta}_{\text{titre \PdeuxOcinq{}}}
\]
Le facteur $\theta$ est le titre. La fonction ne rend pas des tonnes d'acide : elle rend des
tonnes de \PdeuxOcinq{}. \hfill$\square$

\paragraph{Preuve 2 --- par la conservation.}
Développée à la section~\ref{sec:concentration} du chapitre précédent : le rendement vaut
$1$ si et seulement si l'on compte en \PdeuxOcinq{}. \hfill$\square$

\paragraph{Preuve 3 --- par l'industrie.}
EMAPHOS est officiellement dimensionné en \SI{280000}{\tonne}~\PdeuxOcinq{} par
an~\cite{ocp-emaphos}, soit environ \SI{770}{\tonne}~\PdeuxOcinq{} par jour. Sa demande dans
le scénario est de \SI{700}{\tonne} par jour : l'ordre de grandeur ne concorde que dans
cette unité. \hfill$\square$

\begin{alerte}[title={Portée de ce résultat}]
Cette information n'est écrite dans aucun des documents fournis. Sans elle, l'intégralité
des bilans matière du modèle serait erronée. C'est le résultat préliminaire le plus
important de ce travail.

Les rendements de clarification ($0{,}90$) et de cocristallisation ($0{,}80$) s'interprètent
alors comme des \emph{taux de récupération du \PdeuxOcinq{}}, le complément partant dans les
boues --- ce qui est physiquement cohérent et \emph{ferme} le bilan matière global.
\end{alerte}

\section{Nature du problème}

\begin{table}[H]
\centering
\caption{Caractéristiques du modèle (valeurs mesurées sur le scénario de référence)}
\label{tab:nature-modele}
\begin{tabular}{@{}ll@{}}
\toprule
\textbf{Caractéristique} & \textbf{Valeur} \\
\midrule
Type & Programme linéaire mixte en nombres entiers (MILP) \\
Horizon & 1 journée (mono-période) \\
Unité & tonne de \PdeuxOcinq{} \\
Variables continues & 141 \\
Variables binaires & 29 \\
Contraintes & 138 \\
Objectif & lexicographique à 3 niveaux \\
Solveur & PuLP + CBC \\
Temps de résolution & \SI{0.06}{\second} \\
\bottomrule
\end{tabular}
\end{table}

Deux familles de décisions sont intrinsèquement discrètes --- les niveaux de décadmiation et
l'activation des transferts --- ce qui interdit une relaxation continue : celle-ci
produirait des plans physiquement irréalisables, comme décadmier \SI{312}{\tonne}.

\section{Ensembles}

\begin{table}[H]
\centering
\caption{Ensembles d'indices}
\label{tab:ensembles}
\begin{tabular}{@{}llc@{}}
\toprule
\textbf{Symbole} & \textbf{Définition} & \textbf{Cardinal} \\
\midrule
$\Lvingtneuf$ & lignes de production d'acide 29 & 6 \\
$\Lcinqquatre$ & lignes de concentration & 5 \\
$\Ech$ & échelons de concentration & 24 \\
$\Conso$ & consommateurs & 9 \\
$\Prod$ & produits engrais & variable \\
$\Acides$ & types d'acide & 6 \\
\bottomrule
\end{tabular}
\end{table}

\subsection{Partition des types d'acide}

\begin{equation}
\Acides = \underbrace{\{\text{29std},\ \text{29dec}\}}_{\Acides_{29}}
\ \cup\
\underbrace{\{\text{54ncl},\ \text{54cl},\ \text{54coc},\ \text{54deccl}\}}_{\Acides_{54}}
\end{equation}

On distingue par ailleurs deux sous-ensembles selon le mode de livraison :
\begin{align}
\Acides^{\mathrm{loc}} &= \{\text{29std},\ \text{29dec},\ \text{54ncl}\}
&&\text{livrés depuis un stock local} \\
\Acides^{\mathrm{cen}} &= \{\text{54cl},\ \text{54coc},\ \text{54deccl}\}
&&\text{livrés depuis IR11 ou IR12}
\end{align}

\begin{encadre}
\textbf{Justification de cette partition.} Elle traduit la règle « aucun acide traité n'est
expédié directement depuis sa ligne ». Seuls les acides de $\Acides^{\mathrm{loc}}$ sont
soumis aux contraintes d'interconnexion, ce qui divise par deux le nombre de variables de
livraison à créer.
\end{encadre}

\subsection{Application ligne 29 $\to$ ligne 54}

\begin{equation}
\sigma : \Lvingtneuf \longrightarrow \Lcinqquatre \cup \{\varnothing\}
\end{equation}
avec $\sigma(\text{13AB}) = \text{14AB}$, \dots, $\sigma(\text{13E}) = \text{14EXT}$ et
$\sigma(\text{13F}) = \varnothing$.

\begin{proposition}
La restriction de $\sigma$ à $\Lvingtneuf \setminus \{\text{13F}\}$ est une bijection sur
$\Lcinqquatre$.
\end{proposition}

Son inverse $\sigma^{-1}$ est donc bien défini sur $\Lcinqquatre$, ce qui autorise à écrire
indifféremment une équation « côté 29 » ou « côté 54 ». On note
$\Lvingtneufetoile = \Lvingtneuf \setminus \{\text{13F}\}$.

\subsection{Échelons et relations d'incidence}

Les échelons sont partitionnés par ligne, $\Ech = \bigsqcup_{m} \Ech_m$, et chaque ligne
distingue ses échelons cocristallisants $\Ech^{\mathrm{coc}}_m \subseteq \Ech_m$.

Deux relations d'incidence décrivent le réseau : $\Gamma \subseteq \Lvingtneuf \times
\Lvingtneuf$ pour les transferts interzones ($|\Gamma| = 14$, relation anti-réflexive et
\emph{non symétrique}), et $\Phi_a$ pour les raccordements ligne $\to$ consommateur.

\section{Paramètres}

Un paramètre est une donnée connue \emph{avant} la résolution. Le test pratique :
« l'usine peut-elle changer cela demain matin ? » Si non, c'est un paramètre.

\subsection{Grandeurs dérivées de la production}

Calculées en prétraitement à partir des heures de marche :
\begin{align}
\kappa_e &= C_e \cdot \frac{h_e}{24} && \text{production de l'échelon } e \\
\Pi_m &= \sum_{e \in \Ech_m} \kappa_e && \text{production totale de la ligne } m \\
\Picoc_m &= \sum_{e \in \Ech^{\mathrm{coc}}_m} \kappa_e,
\qquad \Pincl_m = \Pi_m - \Picoc_m \\
G_m &= \etacoc \cdot \Picoc_m && \text{CoC produit (systématique)} \\
B^{\mathrm{coc}}_m &= (1 - \etacoc) \cdot \Picoc_m && \text{boue de cocristallisation}
\end{align}

\begin{encadre}[title={Point de méthode essentiel}]
$G_m$ et $B^{\mathrm{coc}}_m$ sont des \textbf{paramètres}, non des variables, parce que la
cocristallisation est systématique. C'est ce qui simplifie le plus le modèle. L'erreur
fréquente consisterait à en faire une décision : le modèle serait plus gros, plus lent, et
faux.
\end{encadre}

\subsection{Demande consolidée}

Un profil qualité est la recette d'un engrais : $q_{p,a}$ tonnes d'acide $a$ par tonne
d'engrais $p$. La conversion demande $\to$ besoins est une application linéaire :
\begin{equation}
\boxed{\;R_{k,a} = \delta_{k,a} + \sum_{p \in \Prod} q_{p,a} \, F_{k,p}\;}
\label{eq:besoins}
\end{equation}
où $\delta_{k,a}$ est la demande directe et $F_{k,p}$ le tonnage d'engrais à produire.

\section{Variables de décision}

\subsection{Décadmiation}

Le choix du niveau est encodé en \emph{one-hot} :
\begin{equation}
w_{\ell,v} \in \{0,1\}, \quad \ell \in \mathcal{L}^{\mathrm{dec}}_{29},\ v \in \mathcal{D}
\qquad\text{et}\qquad
d_\ell = \sum_{v \in \mathcal{D}} v \cdot w_{\ell,v}
\end{equation}

\subsection{Absence de variable pour la clarification décadmiée}

L'acide décadmié concentré est \emph{obligatoirement} clarifié. Plutôt que créer une
variable $\udec_m$ et une contrainte d'égalité, on substitue :
\begin{equation}
\udec_m \equiv \xdec_{\sigma^{-1}(m)}
\label{eq:substitution-udec}
\end{equation}

\begin{encadre}[title={Règle générale de modélisation}]
\emph{Une contrainte d'égalité qui définit une variable doit être éliminée par
substitution.} Ici, cela supprime cinq variables et cinq contraintes, et rend automatique le
couplage entre décadmiation et capacité de décantation.
\end{encadre}

\subsection{Variables d'écart}

Trois familles de variables d'écart complètent le modèle : $\rho_{k,a}$ pour la demande non
satisfaite, $\nu^{+}, \nu^{-}$ pour les dépassements de bandes de stock, et
$\varepsilon^{+}, \varepsilon^{-}$ pour les écarts à la charge de concentration.

\begin{encadre}[title={Pourquoi des variables d'écart}]
Un modèle qui répond « Infeasible » est inutilisable en exploitation : l'ingénieur veut
savoir \emph{de combien} et \emph{où} cela coince, pas se voir refuser une réponse. Les
écarts rendent le modèle toujours faisable, et leur valeur à l'optimum constitue en
elle-même une information de diagnostic.
\end{encadre}

\section{Contraintes}

Les contraintes sont numérotées \Ctr{1} à \Ctr{17}. Les six les plus significatives sont
développées ci-dessous ; l'ensemble figure en annexe~\ref{ann:contraintes}.

\subsection{\Ctr{4} --- Charge de concentration}

\begin{equation}
\boxed{\;\xstd_\ell + \xdec_\ell = \Pi_{\sigma(\ell)} + \varepsilon^{+}_\ell - \varepsilon^{-}_\ell\;}
\qquad \forall \ell \in \Lvingtneufetoile
\label{eq:C4}
\end{equation}
avec $0 \le \varepsilon^{\pm}_\ell \le \epsilon\, \Pi_{\sigma(\ell)}$.

\paragraph{Français.} La ligne 29 fournit à sa concentration ce que celle-ci va produire, à
une tolérance $\epsilon$ près.

\paragraph{Origine.} Interrogé sur cette contrainte, l'encadrant industriel a précisé :
\begin{quote}
\itshape « Oui la somme doit être égale, c'est normalement rigide, mais on peut avoir des
seuils de tolérance pour respecter la planification des transferts. »
\end{quote}

\paragraph{Justification de la forme.} Cette réponse contient deux moitiés apparemment
contradictoires. Une égalité stricte ignorerait la tolérance accordée ; une inégalité libre
la banaliserait, l'optimiseur s'écartant du plan dès que cela l'arrange. La bande
\emph{pénalisée} au niveau 2 de l'objectif rend fidèlement les deux moitiés : le modèle
reste à l'égalité stricte sauf lorsque s'en écarter évite une violation plus grave.

\begin{resultat}[title={Vérification empirique}]
Sur le scénario de référence, la solution optimale n'utilise \textbf{aucune} tolérance : les
cinq lignes sont exactement au plan. Le comportement obtenu est donc bien celui décrit comme
« normalement rigide ».
\end{resultat}

\subsection{\Ctr{4b} --- Alimentation des échelons cocristallisants}

\begin{equation}
\xstd_\ell \ \ge\ \Picoc_{\sigma(\ell)} \qquad \forall \ell \in \Lvingtneufetoile
\end{equation}

Conséquence directe du caractère souple de \Ctr{4} : si la charge peut s'écarter du plan, il
faut préciser \emph{quels} échelons absorbent l'écart. Les unités de cocristallisation ne se
modulant pas, ce sont les échelons ordinaires qui l'absorbent.

\subsection{\Ctr{6} --- NCL réellement disponible}

\begin{equation}
N_m = \xstd_{\sigma^{-1}(m)} - \Picoc_m \qquad \forall m \in \Lcinqquatre
\label{eq:C6}
\end{equation}

\begin{encadre}[title={Cohérence avec la formulation initiale}]
Tant que \Ctr{4} était une égalité stricte, on pouvait écrire $N_m = \Pincl_m - \xdec$. Dès
lors que la charge peut s'écarter du plan, la production doit suivre la charge réelle.

\emph{Vérification} : quand la charge vaut exactement le plan,
$\xstd + \xdec = \Pi_m = \Picoc_m + \Pincl_m$, donc $\xstd - \Picoc_m = \Pincl_m - \xdec$.
La nouvelle écriture \textbf{généralise} l'ancienne sans la contredire --- contrôle à
effectuer systématiquement lorsqu'on relâche une contrainte.
\end{encadre}

\subsection{\Ctr{7} --- Bilan matière du stock d'acide 29 standard}

C'est l'équation la plus riche du modèle : six flux entrants, deux sortants.

\begin{align}
S^{\mathrm{std}}_\ell = \;& S^{0,\mathrm{std}}_\ell + (P_\ell - d_\ell)
&& \text{initial + production} \notag \\
&+ \textstyle\sum_{\ell' : (\ell',\ell) \in \Gamma} t_{\ell'\ell}
 - \textstyle\sum_{\ell' : (\ell,\ell') \in \Gamma} t_{\ell\ell'}
&& \text{transferts} \notag \\
&+ B^{\mathrm{coc}}_{\sigma(\ell)}
&& \text{boue de cocristallisation} \notag \\
&+ (1-\etacl)\left(\uncl_{\sigma(\ell)} + \xdec_\ell\right)
&& \text{boue de clarification} \notag \\
&+ \beta_\ell \textstyle\sum_{m} y^{\text{54ncl}}_{m,\text{EMAPHOS}}
&& \text{retour EMAPHOS} \notag \\
&- \xstd_\ell - \textstyle\sum_{k} y^{\text{29std}}_{\ell,k}
&& \text{concentration + livraisons}
\label{eq:C7}
\end{align}

\begin{alerte}[title={Les trois pièges de cette équation}]
\begin{enumerate}
  \item \textbf{Oublier la boue de clarification de l'acide décadmié.} Le terme
        $(1-\etacl)\,\xdec_\ell$ est facile à manquer, car cette clarification est
        automatique et n'a pas de variable dédiée.
  \item \textbf{Se tromper de niveau pour le retour EMAPHOS.} Il est prélevé au niveau 54 et
        restitué au niveau 29 --- le seul flux du système qui remonte.
  \item \textbf{Inverser les indices des transferts.} $t_{\ell'\ell}$ entre,
        $t_{\ell\ell'}$ sort.
\end{enumerate}
\end{alerte}

\subsection{\Ctr{11} --- Capacité des décanteurs}

\begin{equation}
\boxed{\;\uncl_m + \xdec_{\sigma^{-1}(m)} \ \le\ \Lambda\;}
\qquad \forall m \in \mathcal{L}^{\mathrm{clar}}_{54}
\label{eq:C11}
\end{equation}

C'est la contrainte structurante du modèle, celle qui crée l'arbitrage décrit au
chapitre~\ref{chap:procede}. La limite porte sur le \textbf{débit entrant} : un décanteur
est dimensionné par ce qu'il reçoit, non par ce qu'il restitue.

\subsection{\Ctr{14} --- Satisfaction de la demande}

Trois cas selon la nature de l'acide. Le plus intéressant est celui du besoin en
\emph{acide 54 de qualité décadmiée} :
\begin{equation}
\boxed{\;y^{\mathrm{coc}}_{k} + y^{\mathrm{deccl}}_{k} + \rho_{k,\mathrm{dectot}} = R_{k,\mathrm{dectot}}\;}
\label{eq:C14-dectot}
\end{equation}

\begin{encadre}
Ce besoin ne porte pas sur un \emph{produit} mais sur une \emph{propriété}. Deux produits
physiquement distincts --- le CoC et le DEC\_CL --- la possèdent et sont donc
substituables. La contrainte porte sur leur \textbf{somme}, ce qui exprime exactement cette
substituabilité. C'est aussi pourquoi ils partagent le même bac IR11 : du point de vue du
client, ils sont interchangeables.
\end{encadre}

\subsection{\Ctr{15} --- Contrainte de qualité d'IR11}

Le cahier des charges affirme que la production de DEC\_CL est \emph{obligatoire}, IR11
devant contenir un mélange de CoC et de DEC\_CL sous peine de dégradation de la qualité.
Aucun seuil chiffré n'est toutefois fourni. On modélise cette règle par une contrainte de
composition, d'expression linéaire après réarrangement :
\begin{equation}
(1-\alpha)\,\etacl \sum_{m} \xdec_{\sigma^{-1}(m)} \ \ge\ \alpha \sum_{m} G_m
\label{eq:C15}
\end{equation}

\begin{encadre}[title={Pourquoi cette contrainte est indispensable}]
Sans elle, l'optimiseur ne produirait \emph{jamais} de DEC\_CL : la cocristallisation
systématique (\SI{2431}{\tonne}/j) couvre déjà tout le besoin en acide de qualité décadmiée
(\SI{987}{\tonne}/j), si bien que produire du DEC\_CL est pour lui un coût sans bénéfice.
Il violerait donc la règle métier.

Ce raisonnement a été vérifié expérimentalement : un test automatisé démontre qu'avec
$\alpha = 0$, la production de DEC\_CL tombe exactement à zéro
(chapitre~\ref{chap:implementation}).
\end{encadre}

\subsection{\Ctr{16} --- Transferts semi-continus}

\begin{equation}
\tau^{\min} b_{\ell\ell'} \ \le\ t_{\ell\ell'} \ \le\ \tau^{\max} b_{\ell\ell'}
\qquad \forall (\ell,\ell') \in \Gamma
\label{eq:C16}
\end{equation}

\begin{proof}[Démonstration de correction]
Si $b = 0$, alors $0 \le t \le 0$ donc $t = 0$. Si $b = 1$, alors
$\tau^{\min} \le t \le \tau^{\max}$. Ces deux cas sont exactement ceux souhaités, et aucun
autre n'est atteignable.
\end{proof}

\section{Fonction objectif}

\subsection{L'objectif spécifié est dégénéré}
\label{sec:degenerescence}

Le cahier des charges demande de \emph{maximiser le total d'acide livré}, tout en imposant
que \emph{toutes les demandes soient exactement satisfaites}.

\begin{proposition}[Dégénérescence]
Sous l'hypothèse que la satisfaction de la demande est une contrainte d'égalité, l'objectif
« maximiser le total livré » est constant sur le domaine réalisable.
\end{proposition}

\begin{proof}
Notons $Y$ le total livré. Par la contrainte \Ctr{14} avec $\rho \equiv 0$ :
\[
Y = \sum_{k,a} \big(\text{livraisons vers } k \text{ en acide } a\big) = \sum_{k,a} R_{k,a}
\]
Or $R_{k,a}$ est un paramètre. Donc $Y$ est constant pour toute solution réalisable : la
fonction objectif ne dépend d'aucune variable, et son gradient est nul.
\end{proof}

\begin{alerte}[title={Conséquences pratiques}]
\begin{enumerate}
  \item Le solveur renvoie un sommet \emph{arbitraire} du polyèdre réalisable.
  \item Deux exécutions peuvent produire des plans très différents --- transferts inutiles,
        décadmiation excessive --- avec la \emph{même} valeur d'objectif.
  \item Le résultat perd toute crédibilité auprès des exploitants : « pourquoi le logiciel
        me demande-t-il de transférer \SI{400}{\tonne} aujourd'hui et rien demain, alors que
        rien n'a changé ? »
\end{enumerate}
\end{alerte}

\subsection{La correction : optimisation lexicographique}

On ordonne les objectifs par priorité \emph{stricte}. On optimise $f_1$, on fige son optimum
par une contrainte, on optimise $f_2$ sous cette contrainte, puis $f_3$. Aucun compromis
n'est possible entre niveaux : gagner sur $f_2$ ne justifie jamais de perdre sur $f_1$.

\paragraph{Niveau 1 --- servir les clients.}
\begin{equation}
f_1 = \sum_{k \in \Conso} \sum_{a \in \Acides} \pi_{k,a}\, \rho_{k,a} \ \longrightarrow\ \min
\end{equation}

\begin{encadre}[title={Équivalence avec l'objectif demandé}]
Puisque $\sum_{k,a}(\text{livré}) = \sum_{k,a} R_{k,a} - \sum_{k,a}\rho_{k,a}$ et que le
premier terme est constant, minimiser $f_1$ à poids unitaires revient \emph{exactement} à
maximiser le total livré. Le niveau 1 est donc une généralisation stricte de l'objectif
spécifié.

Il apporte de surcroît une propriété absente de l'original : si le système ne peut pas tout
servir, le modèle ne répond pas « Infeasible » --- il sert au mieux et indique précisément
ce qui manque et pour qui.
\end{encadre}

\paragraph{Niveau 2 --- respecter les bandes de stock.}
\begin{equation}
f_2 = w_{\mathrm{stock}} \sum \left(\nu^{+} + \nu^{-}\right)
    + w_{\mathrm{charge}} \sum \left(\varepsilon^{+} + \varepsilon^{-}\right)
\ \longrightarrow\ \min
\end{equation}

Une cuve qui déborde arrête la ligne amont ; une cuve qui se vide arrête la ligne aval.
Incidents graves, mais moins que de ne pas livrer un client --- d'où leur place au niveau 2.

\paragraph{Niveau 3 --- bien exploiter.}
\begin{equation}
f_3 = c_t \sum t_{\ell\ell'} + c_b \sum b_{\ell\ell'} + c_d \sum d_\ell
    + c_s \sum \left| S - S^{\mathrm{cible}} \right|
\ \longrightarrow\ \min
\end{equation}

\begin{table}[H]
\centering
\caption{Termes du niveau 3 et leur justification industrielle}
\label{tab:niveau3}
\begin{tabular}{@{}lp{8cm}@{}}
\toprule
\textbf{Terme} & \textbf{Justification} \\
\midrule
Volume transféré & énergie de pompage, usure des conduites \\
Nombre de transferts & chaque opération mobilise un opérateur \\
Décadmiation & réactifs, capacité de filtration consommée \\
Écart au milieu de bande & robustesse pour le lendemain \\
\bottomrule
\end{tabular}
\end{table}

\paragraph{Linéarisation de la valeur absolue.}
On pose $S - S^{\mathrm{cible}} = g^{+} - g^{-}$ avec $g^{+}, g^{-} \ge 0$, et l'on minimise
$g^{+} + g^{-}$.

\begin{proof}
Puisqu'on minimise $g^{+} + g^{-}$, l'optimum ne rend jamais les deux simultanément
positifs : on pourrait les diminuer tous deux de la même quantité sans changer leur
différence. À l'optimum, l'un des deux est donc nul et
$g^{+} + g^{-} = |S - S^{\mathrm{cible}}|$.
\end{proof}

\begin{encadre}
C'est le niveau 3 qui donne à la solution sa qualité opératoire : il transforme un plan
simplement réalisable en un plan \textbf{sobre et robuste} --- pas de transfert inutile, pas
de décadmiation superflue, des cuves laissées en bonne position pour le lendemain.
\end{encadre}

\subsection{Mise en œuvre}

\begin{enumerate}
  \item minimiser $f_1$, obtenir $f_1^\star$ ;
  \item ajouter $f_1 \le f_1^\star + \varepsilon_1$, minimiser $f_2$, obtenir $f_2^\star$ ;
  \item ajouter $f_2 \le f_2^\star + \varepsilon_2$, minimiser $f_3$.
\end{enumerate}

Les $\varepsilon_i$ absorbent les erreurs d'arrondi du solveur. Sans eux, figer
$f_1 = f_1^\star$ à l'exact rend l'étape suivante numériquement infaisable : le solveur ne
retrouve pas au bit près la valeur qu'il vient d'annoncer.

\section{Hypothèses du modèle}

Toute hypothèse est un risque assumé. Elle doit être écrite, justifiée et testable.

\begin{table}[H]
\centering
\caption{Hypothèses de modélisation et niveau de risque associé}
\label{tab:hypotheses}
\begin{tabular}{@{}clp{5.8cm}c@{}}
\toprule
\textbf{N\textsuperscript{o}} & \textbf{Hypothèse} & \textbf{Justification} & \textbf{Risque} \\
\midrule
H1 & Unité en tonnes de \PdeuxOcinq{} & trois preuves convergentes & faible \\
H2 & Modèle mono-période & toutes les données sont journalières & faible \\
H3 & Acide dec hors échelons CoC & un CoC ne peut être clarifié & faible \\
H4 & IR11 initial entièrement CoC & hypothèse la plus défavorable & moyen \\
H5 & Boues CoC vers stock 29 std & guide du procédé, référence définitive & moyen \\
H6 & Limite décanteur en entrée & mention « \emph{before yield} » + physique & faible \\
H7 & $\tau^{\max} = \SI{1000}{\tonne}$ & valeur d'attente & \textbf{élevé} \\
H8 & $\alpha = 0{,}15$ & valeur d'attente, sans fondement métier & \textbf{très élevé} \\
H9 & Bornes de stock recalculées & méthode validée sur l'acide 54 & \textbf{élevé} \\
H10 & Bornes hautes molles & sinon le modèle est infaisable & moyen \\
H11 & Rendements déterministes & données de procédé stabilisées & faible \\
H12 & Pas de délai de transport & horizon journalier $\gg$ temps de transfert & faible \\
\bottomrule
\end{tabular}
\end{table}

\begin{encadre}
Les hypothèses H7, H8 et H9 font l'objet d'une analyse de sensibilité au
chapitre~\ref{chap:resultats}. C'est la réponse méthodologique correcte à une donnée
incertaine : ne pas la deviner, mais \textbf{mesurer l'influence de son ignorance}.
\end{encadre}

\section{Traçabilité contrainte $\leftrightarrow$ réalité industrielle}

\begin{table}[H]
\centering
\caption{Correspondance entre le discours d'atelier et sa traduction mathématique}
\label{tab:tracabilite}
\begin{tabular}{@{}p{7.1cm}p{6.2cm}@{}}
\toprule
\textbf{Réalité de l'atelier} & \textbf{Traduction mathématique} \\
\midrule
« La production d'acide 29 nous est imposée » & $P_\ell$ paramètre \\
« On allume 1 ou 2 filtres, pas 1,4 » & $d_\ell \in \{0,750,1500\}$, binaires \\
« Concentrer conserve le \PdeuxOcinq{} » & $\etaconc = 1$ \\
« Les échelons tournent selon le planning » & $\Pi_m$ paramètre, \Ctr{4} \\
« La cocristallisation traite tout, tout le temps » & $G_m$ paramètre, aucune variable \\
« L'acide décadmié est clarifié immédiatement » & substitution~\eqref{eq:substitution-udec} \\
« Les 2 décanteurs font \SI{1000}{\tonne}/jour » & \Ctr{11}, équation~\eqref{eq:C11} \\
« Les boues, on les recycle » & termes de retour dans~\eqref{eq:C7} \\
« EMAPHOS nous renvoie \SI{40}{\percent} » & terme $\beta_\ell$ dans~\eqref{eq:C7} \\
« Toutes les lignes ne sont pas reliées » & $\Gamma$, $\Phi_a$ \\
« On ne transfère pas 4 tonnes » & \Ctr{16}, semi-continue \\
« La cuve ne doit ni déborder ni se vider » & \Ctr{9}, \Ctr{17}, molles \\
« IR11 doit contenir un mélange » & \Ctr{15}, équation~\eqref{eq:C15} \\
« Le client d'abord » & niveau 1 de l'objectif \\
« Sans gaspiller » & niveau 3 de l'objectif \\
\bottomrule
\end{tabular}
\end{table}
